\documentclass[10pt,a4paper]{amsart}
\usepackage[margin=19mm]{geometry}
\usepackage{amsmath,amssymb,mathtools}
\usepackage[scr=rsfs,cal=euler]{mathalfa}
\usepackage{xfrac}
\usepackage[unicode,hidelinks,bookmarks=false]{hyperref}
\newcommand{\ie}{i.e.\ }
\newcommand{\cf}{cf.\ }
\newcommand{\resp}{respectively\ }
\newcommand{\Subsection}{\subsection}
\providecommand{\nobreakdash}{\nobreak\hbox{-}}
\setlength{\emergencystretch}{2em}
\multlinegap0pt
\usepackage{bm}
\usepackage{bbm}
\usepackage{dsfont}
\usepackage{xspace}
\usepackage{cancel}
\usepackage{mathtools}
\usepackage{amsthm}
\makeatletter
\@ifundefined{theorem}{\newtheorem{theorem}{Theorem}}{}
\@ifundefined{lemma}{\newtheorem{lemma}[theorem]{Lemma}}{}
\@ifundefined{proposition}{\newtheorem{proposition}[theorem]{Proposition}}{}
\makeatother
\newcommand{\EDS}{\operatorname{EDS}}
\newcommand{\F}{\mathbb{F}}
\newcommand{\Int}{\operatorname{int}}
\newcommand{\Z}{\mathbb{Z}}
\newcommand{\id}{\operatorname{id}}
\title[Unbounded Gaps Between Ordinary and Equivariant Dehn Surgery]{Unbounded Gaps Between Ordinary and Equivariant Dehn Surgery Numbers}
\author{Qilong Guo, Chunxing Yan}
\date{Source: arXiv:2608.03886v1 (CC BY 4.0)}
\begin{document}
\maketitle

\noindent\emph{Source and licence.} Unbounded Gaps Between Ordinary and Equivariant Dehn Surgery Numbers. Version arXiv:2608.03886v1 (CC BY 4.0), distributed under the Creative Commons Attribution 4.0 International licence, as recorded in the arXiv licence metadata for this exact version and shown on the abstract page https://arxiv.org/abs/2608.03886v1. Full rights notice: licenses/arxiv-2608.03886-CC-BY-4.0.txt.\par\smallskip

\noindent\textit{Editorial scope.} Counted unit: the Proposition whose source label is \texttt{prop:homological-bound} in the article (source0.tex lines 397--402, source label \texttt{prop:homological-bound}), delivered together with the article's own complete original proof of it (source0.tex lines 404--449, one \texttt{proof} environment of 46 lines). No mathematics is written, repaired or re-derived: statement and proof are the article's own text line for line. Required context retained uncounted: lem:filling (lines 219--225); lem:signed-permutation (lines 380--394). Every definition, display and label of those lines is retained unchanged. Omitted from this package and not redistributed: 1--218, 226--379, 395--396, 403--403, 450--685. Nothing inside a retained excerpt is omitted or altered: every retained body is delivered exactly as the article's lines are written, so no omission receipt of any kind accompanies this package. Citations: the retained excerpts carry no citation command at all, so no bibliography entry of the article is transcribed and no reference is redistributed. Reference adaptation: none was needed and none was made; every reference inside the delivered unit resolves inside it, so no unresolved reference or citation is produced. Printed slips kept literal: no printed word, symbol, hypothesis or number of the counted statement or of its proof was altered. Layout and class: the article's own class is \texttt{amsart} with packages including fontenc, lmodern, microtype, amsmath, amssymb, mathtools, hyperref; the wrapper is the lane's pinned \texttt{amsart} editorial wrapper at 10pt on A4, which restates the theorem-like environments the retained text needs and adds the article's own macro definitions that the retained text needs (\texttt{\textbackslash EDS}, \texttt{\textbackslash F}, \texttt{\textbackslash Int}, \texttt{\textbackslash Z}, \texttt{\textbackslash id}), with extra wrapper packages (bm, bbm, dsfont, xspace, cancel, mathtools). The delivered numbering is the unit's own. Assets: no retained span contains a figure environment, a graphics-inclusion command, a PDF-page inclusion, a table or a TikZ picture; no image file is copied into this package, the package declares no asset dependency and no image file is redistributed with it. Licence: version arXiv:2608.03886v1 of this article is distributed under the Creative Commons Attribution 4.0 International licence, as recorded in the arXiv licence metadata for this exact version and shown on the abstract page https://arxiv.org/abs/2608.03886v1. A full rights notice is delivered at licenses/arxiv-2608.03886-CC-BY-4.0.txt. Characters: every retained excerpt is pure ASCII, so the delivered engine, XeLaTeX with its default Latin Modern text font, reports no missing character.
\begin{lemma}\label{lem:filling}
Let $X$ be the exterior of an $m$-component link in $S^3$, and let $Y$ be obtained by filling every boundary component of $X$. For every commutative coefficient ring $R$, inclusion induces a surjection
\[
   H_1(X;R)\twoheadrightarrow H_1(Y;R).
\]
When the filling is equivariant, the surjection is equivariant.
\end{lemma}

\begin{lemma}\label{lem:signed-permutation}
Let $V\cong\Z^m$ have basis $e_1,\dots,e_m$, and let $T\colon V\to V$ be an involution satisfying
\[
   T(e_i)=\varepsilon_i e_{\pi(i)},
   \qquad
   \varepsilon_i\in\{\pm1\}.
\]
There are nonnegative integers $r_+$ and $r_-$ with $r_++r_-=m$ such that, over every field $\Bbbk$ of characteristic different from $2$,
\[
   \dim_{\Bbbk}\ker(T-\id)=r_+,
   \qquad
   \dim_{\Bbbk}\ker(T+\id)=r_-.
\]
In particular, these dimensions do not depend on the field.
\end{lemma}

\begin{proposition}\label{prop:homological-bound}
For every closed oriented $3$-manifold $Y$ with an orientation-preserving involution $\tau$,
\[
   \EDS(Y,\tau)\geq b^+(Y,\tau)+b^-(Y,\tau).
\]
\end{proposition}

\begin{proof}
Let $L$ be an $m$-component integral periodic surgery description inducing $(Y,\tau)$, and let
\[
   X=S^3\setminus\Int N(L)
\]
be the link exterior. Choose orientations on the components of $L$. Their meridians form a basis of $H_1(X;\Z)\cong\Z^m$. The ambient half-turn permutes the components and sends each oriented meridian to a meridian or its negative. Its action $T$ on $H_1(X;\Z)$ is therefore a signed permutation involution.

By Lemma~\ref{lem:signed-permutation}, there are integers $r_+,r_-\geq0$, independent of the odd prime $\ell$, such that
\[
   \begin{aligned}
   H_1(X;\F_\ell)&=W^+_\ell\oplus W^-_\ell,\\
   \dim W^+_\ell&=r_+,
   \qquad
   \dim W^-_\ell=r_-,
   \qquad
   r_++r_-=m.
   \end{aligned}
\]
Lemma~\ref{lem:filling} gives an equivariant surjection
\[
   f_\ell\colon H_1(X;\F_\ell)\twoheadrightarrow H_1(Y;\F_\ell),
   \qquad
   f_\ell T=\tau_*f_\ell.
\]
The restriction of $f_\ell$ to each eigenspace is surjective. To see this, let $y$ be a $\delta$-eigenvector and choose $x$ with $f_\ell(x)=y$. Then
\[
   x_\delta=\frac{x+\delta T(x)}2
\]
is a $\delta$-eigenvector and satisfies $f_\ell(x_\delta)=y$. It follows that
\[
   b^+_\ell(Y,\tau)\leq r_+,
   \qquad
   b^-_\ell(Y,\tau)\leq r_-
\]
for every odd prime $\ell$. Since $r_+$ and $r_-$ are independent of $\ell$, the two maxima may be taken separately:
\[
   b^+(Y,\tau)\leq r_+,
   \qquad
   b^-(Y,\tau)\leq r_-.
\]
It follows that
\[
   b^+(Y,\tau)+b^-(Y,\tau)\leq r_++r_-=m.
\]
Minimizing over all periodic surgery descriptions proves the proposition.
\end{proof}

\end{document}
